%%%%%%%%%%%%%%%%%%%%%%%%%%%%%%%%%%%%%%%%%%%%%%%%%%%%%%%%%%%%%%%%%%%%%%
%                                                                    %
% Uniform two-weight norm inequalities for Hankel transform          %
% Bochner-Riesz means of order one,                                  %
% by                                                                 %
% \'Oscar Ciaurri, Krzysztof Stempak and Juan L. Varona              %
%                                                                    %
% DATE: 6-November-2003                                              %
%                                                                    %
% Written with AMS-LaTeX 2.0                                         %
%                                                                    %
%%%%%%%%%%%%%%%%%%%%%%%%%%%%%%%%%%%%%%%%%%%%%%%%%%%%%%%%%%%%%%%%%%%%%%

\documentclass[leqno,a4paper]{amsart}

\newtheorem{Prop}{Proposition}
\newtheorem{Lem}{Lemma}
\newtheorem{Theorem}{Theorem}
%% \newtheorem{Cor}{Corollary} % UNUSED
\newtheorem*{TheoremKnown}{Theorem}

% \usepackage[notcite,notref]{showkeys}

\newcommand{\N}{\mathbb{N}}
\newcommand{\C}{\mathbb{C}}
\newcommand{\R}{\mathbb{R}}
\newcommand{\Diag}{\mathop{\rm Diag}\nolimits}
\newcommand{\Var}{\mathop{\rm Var}\nolimits}
\newcommand{\overnu}{\overline{\nu}}

%%%%%%%% A LOT OF BIG SEVERAL-LINE EQUATIONS (aligns, etc.).
%%%%%%%% THUS, WE ALLOW TO BREAK THEM IN PAGES.
% \allowdisplaybreaks

%%%%%%%%%%%%%%%%%%%%%%%%%%%%%%%%%%%%%%%%%%%%%%%%%%%%%%
\begin{document}
%%%%%%%%%%%%%%%%%%%%%%%%%%%%%%%%%%%%%%%%%%%%%%%%%%%%%%

\title[Uniform two-weight norm inequalities]{Uniform two-weight
norm inequalities for\\ Hankel transform Bochner-Riesz means\\
of order one}

\author[\'O. Ciaurri]{\'Oscar Ciaurri}
\address{Departamento de Matem\'aticas y Computaci\'on,
Universidad de La Rioja,
% Edificio J.~L.~Vi\-ves,
Calle Luis de Ulloa s/n,
26004 Logro\~no, Spain}
\email{oscar.ciaurri@dmc.unirioja.es}

\thanks{Research of the first and third authors supported by grant
BFM2003-06335-C03-03 of the~DGI}

\author[K. Stempak]{Krzysztof Stempak}
\address{Instytut Matematyki,
Politechnika Wroc\l{a}wska,
Wyb.~Wyspia\'nskiego 27,
50-370 Wro\-c\l{a}w,
Poland}
\email{stempak@im.pwr.wroc.pl}

\author[J. L. Varona]{Juan L. Varona}

\address{Departamento de Matem\'aticas y Computaci\'on,
Universidad de La Rioja,
% Edificio J.~L.~Vi\-ves,
Calle Luis de Ulloa s/n,
26004 Logro\~no, Spain}
\email{jvarona@dmc.unirioja.es}

\subjclass[2000]{Primary 42C10; Secondary 44A20}
%\subjclass{Primary 42C10; Secondary 44A20}

\keywords{Bochner-Riesz means, Hankel transform, Bessel functions, two-weight norm inequalities}

\thanks{Research of the second author supported in part by KBN grant
\#~2 P03A 034 20}

\thanks{THIS PAPER HAS BEEN PUBLISHED IN
\textit{Tohoku Math.\ J. (2)} \textbf{56} (2004), no.~3, 371--392.}

\begin{abstract}
Two-weight $L^p$ norm inequalities, uniform with respect to the order of the
involved Bessel function, are proved for the Bochner-Riesz means of the first
order for the Hankel transform. Both sufficient and necessary conditions for
parameters used in the two weights are determined. The proof relies on uniform
pointwise asymptotic estimates for the Bessel functions that were shown by
Barcel\'o and C\'ordoba.
\end{abstract}

%%%%%%%%%%%%%%%%%%%%%%%%%%%%%%%%%%%%%%%%%%%%%%%%%%%%%%
\maketitle
%%%%%%%%%%%%%%%%%%%%%%%%%%%%%%%%%%%%%%%%%%%%%%%%%%%%%%

%%%%%%%%%%%%%%%%%%%%%%%%%%%%%%%%%%%%%%%%%%%%%%%%%%%%%%
\section{Introduction}
%%%%%%%%%%%%%%%%%%%%%%%%%%%%%%%%%%%%%%%%%%%%%%%%%%%%%%
The Hankel transform ${\mathcal H}_\nu f$ of a suitable function $f$ on
$(0,\infty)$ is defined by
$$
{\mathcal H}_\nu f(x)
= \int_0^\infty (xy)^{1/2}J_\nu(xy)f(y)\,dy,\qquad x>0.
$$
Here $\nu>-1$ is given and $J_\nu(x)$ denotes the Bessel function of the
first kind and order $\nu$, \cite{Le}. It is known that
$({\mathcal H}_\nu\circ{\mathcal H}_\nu)f=f$ and $\|{\mathcal H}_\nu
f\|_2=\|f\|_2$, for any
$f\in C_c^\infty(0,\infty)$, the space of $C^\infty$ functions with
compact support in $(0,\infty)$ ($\|\cdot\|_p$ denotes the usual unweighted
norm in $L^p(0,\infty)$). Consider the Bochner-Riesz means
of order $\delta\ge0$ for the Hankel transform ${\mathcal H}_\nu$ given by
\begin{align*}
S_{\nu,R}^\delta f(x)
&= {\mathcal H}_\nu(m_R^\delta
\cdot{\mathcal H}_\nu f)(x)\\
&= \int_0^\infty K_{\nu,R}^\delta(x,y)f(y)\,dy,
\end{align*}
where $m_R^\delta(y)=(1-(y/R)^2)^\delta$ for $0<y<R$ and 0 otherwise, and
\begin{equation*}
K^{\delta}_{\nu,R}(x,y) = \int_0^R m_R^\delta(s)
(xs)^{1/2}J_\nu(xs)(ys)^{1/2}J_\nu(ys)\, ds.
\end{equation*}
Uniform boundedness, with respect to $R>0$, of $S_{\nu,R}^\delta$ in weighted
$L^{p}$ spaces for fixed $\nu>-1$ and $\delta=0$,
\begin{equation}
\|S_{\nu,R}^0 f(x)x^a\|_{p}\le C\|f(x)x^a\|_{p},\qquad R>0,
\label{1a}
\end{equation}
is known to hold if either $\nu\ge-1/2$, $1<p<\infty$ and $-1/p<a<1-1/p$
or $-1<\nu<-1/2$, $2/(2\nu+3)<p<-2/(2\nu+1)$ and $a=0$; cf.~\cite{St}
for appropriate references. The general case $\delta>0$ then
follows by applying a sort of Ste\v ckin-type multiplier theorem for the Hankel
transform. Indeed, let $m(x)$ be a function of bounded variation
on $(0,\infty)$. Then (\ref{1a}) implies
$$
\|{\mathcal H}_\nu(m\cdot{\mathcal H}_\nu f)(x)x^a\|_p\le C\max\{|m(1)|,
\Var m\}\|f(x)x^a\|_{p}
$$
(with the same constant $C$). To show this, it is sufficient to establish
the inequality
$$
\bigg|\int_0^\infty m(y){\mathcal H}_\nu f(y){\mathcal H}_\nu g(y)\,dy\bigg|
\le C\max\{|m(1)|,\Var m\}\|f(x)x^a\|_p\|g(x)x^{-a}\|_{p'}
$$
for $f,g\in C_c^\infty(0,\infty)$. Define
$$
n_\nu(R)=\int_0^R{\mathcal H}_\nu f(y){\mathcal H}_\nu g(y)\,dy, \qquad R>0.
$$
Then (\ref{1a}) gives
$$
|n_\nu(R)|\le C\|f(x)x^a\|_p\|g(x)x^{-a}\|_{p'}
$$
and an integration by parts leads to
\begin{align*}
\bigg|\int_\varepsilon^N m(y){\mathcal H}_\nu f(y){\mathcal H}_\nu g(y)\,dy\bigg|
&\le |m(\varepsilon)n_\nu(\varepsilon)|+|m(N)n_\nu(N)|
+\bigg|\int_\varepsilon^N n_\nu(y)\,dm(y)\bigg|\\
&\le C\max\{|m(1)|,\Var m\}\|f(x)x^a\|_p\|g(x)x^{-a}\|_{p'}.
\end{align*}
Finally, a limit passage with $\varepsilon$ tending to 0 and $N$ tending to
$\infty$ completes the argument.

Since (\ref{1a}), with $\delta=0$ and $R=1$ (hence also with any $R>0$),
holds uniformly with respect to $\nu\ge2$ (cf.~\cite{St} for appropriate
comments) in the case when $4/3<p<4$ and $-1/p<a<1-1/p$ the same remains valid,
by applying the argument just used, for arbitrary $\delta>0$.

The present paper focuses on proving more general, uniform with respect to
$\nu$, inequalities of the form
$$
\|S_\nu^1 f(x)x^a(1+x)^{b-a}\|_p\le C\|f(x)x^A(1+x)^{B-A}\|_p;
$$
we simplify the notation by writing $S_\nu^1$ and $K_\nu^1$
in place of $S_{\nu,1}^1$ and~$K_{\nu,1}^1$.

%%%%%%%%%%%%%%%%%%%%%%%%%%%%%%%%
\begin{Theorem}
\label{Theorem:main}
Let $\alpha>-1$ and $1\le p\le\infty$. Then
\begin{equation}
\|S_\nu^1 f(x)x^a(1+x)^{b-a}\|_p\le C\|f(x)x^A(1+x)^{B-A}\|_p
\label{t4}
\end{equation}
with a constant $C=C(p,\alpha,a,b,A,B)$ independent of $\nu\ge\alpha$ and~$f$
if and only if $a,b,A,B, \alpha$ and $p$ satisfy the conditions
\begin{equation}
a>-1/p-(\alpha+1/2)\,\,\,(\ge \text{if } p=\infty), \qquad
A<1-1/p+(\alpha+1/2)\,\,\,(\le \text{if } p=1),
\label{t1}
\end{equation}
\begin{equation}
b<2-1/p\,\,\,(\le \text{if } p=\infty),
\qquad  B>-1-1/p\,\,\,(\ge \text{if }p=1),
\label{t2}
\end{equation}
\begin{equation}
b \le B.
\label{t3}
\end{equation}
Moreover, only $\nu=\alpha$ in $(\ref{t4})$ is required to prove necessity of
$(\ref{t1})$, $(\ref{t2})$ and $(\ref{t3})$.
\end{Theorem}
%%%%%%%%%%%%%%%%%%%%%%%%%%%%%%%%

It is tacitly understood that when assuming (\ref{t4}) to hold and then proving the
necessity of (\ref{t1}), (\ref{t2}) and (\ref{t3}) only such functions
$f$ are admitted for which $S^1_\alpha f(x)$ is well defined for (almost) all
$x>0$, i.e.,\ $\int_0^\infty|K^1_\alpha(x,y)f(y)|\,dy < \infty$ for $x>0$.
Clearly, if $f$ is compactly supported in $(0,\infty)$, say in $[c,d]$, $0<c<d<\infty$,
and $\int_c^d|f(y)|\,dy < \infty$, then $\int_0^\infty|K^1_\nu(x,y)f(y)|\,dy<\infty$,
since, with given $x>0$, $K^1_\nu(x,y)$ is a continuous function of $y>0$.
But if we assume $\int_0^\infty|f(x)x^A(1+x)^{B-A}|^p\,dx < \infty$
with $A<1-1/p+(\alpha+1/2)$ ($\le$ if $p=1$) and $B>-1-1/p$ ($\ge$ if $p=1$),
then $S_\nu^1 f(x)$ is also well-defined for $\nu\ge\alpha$ and $x>0$.
This easily follows, by using H\"older's inequality and the estimates
$K_\nu^1(x,y) = O(y^{\nu+1/2})$, $y\to0^+$, and $K_\nu^1(x,y)=O(y^{-2})$,
$y\to\infty$. The first estimate follows from the integral representation of
$K_\nu^1(x,y)$ and the asymptotic
\begin{equation}
J_\nu(t) = O(t^{\nu}), \qquad t\to0^+,
\label{asy}
\end{equation}
the second is a consequence of the Campbell-type representation
of $K_\nu^1(x,y)$:
cf.~(\ref{ec:jk}) and follow the argument in the proof of Lemma~\ref{Lem:L2}
(now, with $x$ and $\nu$ fixed!).
Therefore the second parts of (\ref{t1}) and (\ref{t2}) should be considered
as ``natural'' assumptions. An analysis also shows that the first conditions
in (\ref{t1}) and (\ref{t2}) are dual to their right counterparts.

It is instructive to compare the assumptions (\ref{t1}), (\ref{t2})
and (\ref{t4}) of Theorem~\ref{Theorem:main} with those of Theorem 1.1 of
\cite{St} (partial sums case). The conditions on $a$ and $A$ are exactly the
same; this is probably explained by the fact that the kernels $K^1_\nu(x,y)$
and $K^0_\nu(x,y)$ have the same behaviour in the neighbourhood of $(0,0)$.
The conditions on $b$ and $B$ are now relaxed (by one from both sides) when
compared with those from Theorem 1.1 of \cite{St}. The most important difference
is with the condition on the relation between $b$ and $B$; it becomes now
very simple when compared with that of Theorem 1.1 of \cite{St}. Furthermore,
the uniformity of (\ref{t4}) for large $\nu$ does not imply additional restrictions:
the condition $b\le B$ is sufficient for (\ref{t4}) to hold for $\nu\ge\alpha$
as well as for the single~$\nu=\alpha$.

In \cite{St} we compared a similar result proved there for partial sum operators
(the case of $\delta=0$) with Muckenhoupt's result \cite{Mu} proved for partial
sum operators for expansions with respect to the system of
Laguerre functions
$$
\psi_n^\alpha(x)=\bigg(\frac{2\Gamma(n+1)}
{\Gamma(n+\alpha+1)}\bigg)^{1/2}\,e^{-x^2/2}\,x^{\alpha+1/2}L_n^\alpha(x^2),
$$
which form a complete orthonormal system in $L^2(0,\infty)$ (for a motivation
leading to such a comparition see again \cite{St}).

It seems worthy to compare our present result with that of Muckenhoupt and Webb
\cite{Mu-We} specified to the first order Ces\`aro means and reformulated
for expansions with respect to~$\{\psi^\alpha_n\}_{n=0}^\infty$,
$$
\sigma_N^{(\alpha,1)}(f,x) = \frac1{N+1}\sum_{n=0}^N(N+1-n)
\langle f,\psi^\alpha_n\rangle\psi^\alpha_n(x).
$$
An earlier result of Poiani,~\cite{Po}, admitted the case $a=A$ and $b=B$
only and did not include some possibilities.

%%%%%%%%%%%%%%%%%%
\begin{TheoremKnown}[\cite{Mu-We}, Theorem~(2.29)]
Let $\alpha>-1$ and $1\le p\le\infty$.
Then
\begin{equation}
\sup_{N\ge0}\|\sigma_N^{(\alpha,1)}(f,x)x^a(1+x)^{b-a}\|_p
\le C\|f(x)x^A(1+x)^{B-A}\|_p,
\label{m6}
\end{equation}
holds with $C=C(p,\alpha,a,b,A,B)$ independent of $f$ if and only if
$a,b,A,B$ and $\alpha$ satisfy the conditions
\begin{equation}
a\ge-1-1/p, \qquad a>-1/p-(\alpha+1/2)\,\,\,(\ge \text{if } p=\infty),
\label{m1}
\end{equation}
%%%%%%
\begin{equation}
A\le a,\qquad A\le 1-1/p+1/2,\qquad A<1-1/p+(\alpha+1/2)
\,\,\,(\le \text{if } p=1),
\label{m2}
\end{equation}
%%%%%%
\begin{equation}
a+B\ge\min\{-3-2/(3p),-2-2/p\},\qquad A+b\le\min\{4-2/p,11/3-2/(3p)\},
\label{m3}
\end{equation}
%%%%%%
\begin{equation}
b\le\min\{3-1/p,8/3+1/(3p)\},\qquad B\ge\max\{-3+1/(3p),-2-1/p\},
\label{m4}
\end{equation}
%%%%%%
\begin{equation}
b\le B+\min\{0,3-4/(3p),5/3+4/(3p)\}.
\label{m5}
\end{equation}
\end{TheoremKnown}
%%%%%%%%%%%%%%%%%%

The last parts of the assumptions (\ref{m1}) and (\ref{m2}) are identical
with our assumptions on $a$ and $A$, cf.~(\ref{t1}). This is caused by the
fact that the kernel $K^1_\alpha(x,y)$ and the kernels that correspond to
$\sigma^{(\alpha,1)}_N$ have the same behaviour for small $x$ and $y$. To see
this, one has to compare the result of Lemma~\ref{Lem:L1} below with the estimate
(2.3) of \cite{Mu-We} (because of the reformulation mentioned above, in
(2.3) of \cite{Mu-We} one has to take $x^2, y^2$ in place of $x, y$, and then
multiply both sides by $(xy)^{1/2}$). In general, however, the assumptions
(\ref{m1})--(\ref{m5}) from the Laguerre case are much more involved than
those from Theorem~\ref{Theorem:main}.
This is again explained by more complicated nature of estimates of the
first order Ces\`aro kernels in the Laguerre case (cf.~(2.3), (2.4)
and (2.2) in \cite{Mu-We} with necessary modifications) when compared
with the estimates contained in Lemmas~\ref{Lem:L1}--\ref{Lem:L4} below.

%%%%%%%%%%%%%%%%%%%%%%%%%%%%%%%%%%%%%%%%%%%%%%%%%%%%%%
\section{Preliminaries}
%%%%%%%%%%%%%%%%%%%%%%%%%%%%%%%%%%%%%%%%%%%%%%%%%%%%%%

We start with writing an exact expression for the kernel
$$
K^{1}_{\nu}(x,y) = \sqrt{xy}\int_0^1 s(1-s^2)J_\nu(xs)J_\nu(ys)\, ds
$$
that corresponds to the Bochner-Riesz mean $S_\nu^1$. We use the notation
$$
u=x^2-y^2,\qquad\qquad v=x^2+y^2
$$
and
\begin{align*}
F_1(x,y) &= \sqrt{xy}J_\nu(x)J_\nu(y),\qquad\qquad
F_2(x,y) = \sqrt{xy}J'_\nu(x)J_\nu(y),\\
F_3(x,y) &= \sqrt{xy}J_\nu(x)J'_\nu(y),\qquad\qquad
F_4(x,y) = \sqrt{xy}J'_\nu(x)J'_\nu(y).
\end{align*}
Then we have
\begin{align}
\label{Cam}
&K_\nu^1(x,y) \\
\nonumber
&\quad= \frac{2}{u^{2}} \left((2\nu^2-v)F_1(x,y)-2x\frac vuF_2(x,y)
+ 2y\frac vuF_3(x,y)-2xyF_4(x,y)\right)
\end{align}
whenever $x\neq y$. This may be checked by noting that
\begin{equation}
\label{dif}
K^1_\nu(x,y) = 2\mathcal{D}^{2}F_1(x,y)
\end{equation}
where
\begin{equation*}
\mathcal{D} = \frac{1}{u}\left(y\frac{\partial}{\partial y}
-x\frac{\partial}{\partial x}\right),
\end{equation*}
using $\mathcal{D}(fgh) = \mathcal{D}(f)gh+f\mathcal{D}(g)h+fg\mathcal{D}(h)$
and $\mathcal{D}(x^ay^bu^{-c})=x^ay^bu^{-c-2}((b-a)u+2cv)$, $a,b,c\in\R$,
and observing how $\mathcal{D}$ acts on $F_i$, $i=1,2,3$:
\begin{align*}
\mathcal{D}(F_1) &= \frac1 u(-x F_2+ y F_3),\\
\mathcal{D}(F_2) &= \frac1u\left(-\frac{\nu^2-x^2}{x}F_1+ F_2+ y F_4\right),\\
\mathcal{D}(F_3) &= \frac1u\left(\frac{\nu^2-y^2}{y}F_1- F_3- x F_4\right);
\end{align*}
the identities are obtained by using the fact that $J_\nu(x)$ satisfies
Bessel's differential equation $J''_\nu(x) = ((\nu^2-x^2)/x^2)J_\nu(x)-
(1/x)J_\nu'(x)$. To verify (\ref{dif}) we first use Lommel's formula
\begin{equation*}
\int_0^stJ_{\nu}(xt)J_{\nu}(yt)\,dt
= s\frac{xJ_{\nu+1}(sx)J_{\nu}(sy)-yJ_{\nu}(sx)J_{\nu+1}(sy)}{x^2-y^2},
\end{equation*}
with $s=1$, the identity $xJ_{\nu+1}(x)=\nu J_{\nu}(x)-xJ'_{\nu}(x)$
and the expression on $\mathcal{D}(F_1)$
to obtain $K^0_\nu(x,y)=\mathcal{D}F_1(x,y)$, where $K^0_\nu(x,y)$ denotes
the integral kernel of $S^0_\nu=S^0_{\nu,1}$.
Then, integrating by parts, using Lommel's formula and
the expression on $\mathcal{D}(F_1)$ shows that
\begin{align*}
K^{1}_\nu(x,y)&=\sqrt{xy}\int_0^1(1-s^2)sJ_{\nu}(xs)J_\nu(ys)\, ds\\
&=2\sqrt{xy}\int_0^1 s\cdot
s\frac{xJ_{\nu+1}(sx)J_{\nu}(sy)-yJ_{\nu}(sx)J_{\nu+1}(sy)}{x^2-y^2}\, ds\\
&=2\sqrt{xy}\int_0^1 s\mathcal{D}(J_\nu(sx)J_\nu(sy))\,ds\\
&=2\mathcal{D}K^{0}_\nu(x,y).
\end{align*}
This completes checking (\ref{dif}). We should add at this point that according
to the Laguerre case, (\ref{dif}) could be called a Campbell-type formula.

If $\nu=-1/2$ then $J_{-1/2}(u) = (2/(\pi u))^{1/2}\cos u$ and ${\mathcal H}_{-1/2}$
reduces to the cosine transform. A direct calculation of the integral
representing $K^1_{-1/2}(x,y)$ then shows that
\begin{align*}
K^1_{-1/2}(x,y) = \frac2\pi\bigg(\frac1{(x-y)^3}\sin(x-y)
&-\frac1{(x-y)^2}\cos(x-y) + \frac1{(x+y)^3}\sin(x+y)\\
&-\frac1{(x+y)^2}\cos(x+y)\bigg).
\end{align*}
Similarly, if $\nu=1/2$ then $J_{1/2}(u) = (2/(\pi u))^{1/2}\sin u$,
${\mathcal H}_{1/2}$ reduces to the sine transform and
\begin{align*}
K^1_{1/2}(x,y) = \frac2\pi\bigg(\frac1{(x-y)^3}\sin(x-y)
&-\frac1{(x-y)^2}\cos(x-y) - \frac1{(x+y)^3}\sin(x+y)\\
&+\frac1{(x+y)^2}\cos(x+y)\bigg).
\end{align*}
These two exact formulas (consistent with (\ref{Cam})!) give an
immediate insight on how estimates of
$K^1_\nu(x,y)$ should look like, for a single $\nu$, in different regions
of the quarter plane. For estimates uniform on $\nu$, (\ref{Cam}) together
with the uniform estimates that follow are required.

We will make an extensive use of the following uniform
pointwise estimates of Bessel functions and their first derivatives: with a
constant $D$,
independent of $\nu$ and $x$, for $\nu\ge1$,
\begin{equation}
|x^{1/2}J_\nu(x)|\le D
\begin{cases}
x^{\nu+1/2}(\Gamma(\nu+1)2^\nu)^{-1},& 0<x<\nu/2,\\
\nu^{1/4}(\nu^{1/3}+|x-\nu|)^{-1/4},& \nu/2<x<2\nu,\\
1,& 2\nu<x<\infty,
\end{cases}
\label{1b}
\end{equation}
and
\begin{equation}
|x^{1/2}J_\nu'(x)|\le D
\begin{cases}
\nu x^{\nu-1/2}(\Gamma(\nu+1)2^\nu)^{-1},& 0<x<\nu/2,\\
\nu^{-1/4}(\nu^{1/3}+|x-\nu|)^{1/4},& \nu/2<x<2\nu,\\
1,& 2\nu<x<\infty,
\end{cases}
\label{2b}
\end{equation}
and, for $-1<\nu<1$,
\begin{equation}
|x^{1/2}J_\nu(x)|\le D
\begin{cases}
x^{\nu+1/2}(\Gamma(\nu+1))^{-1},& 0<x<1/2,\\
1,& 1/2<x<\infty,
\end{cases}
\label{3b}
\end{equation}
and
\begin{equation}
|x^{1/2}J_\nu'(x)|\le D
\begin{cases}
x^{\nu-1/2}(\Gamma(\nu+1))^{-1},& 0<x<1/2,\\
1,& 1/2<x<\infty.
\end{cases}
\label{4b}
\end{equation}
In the case $\nu=0$, the bound $x^{-1/2}$ in (\ref{4b}) has to be replaced
by $x^{3/2}$. (\ref{1b}) and (\ref{2b}) are direct consequences of the delicate
bounds done by Barcel\'o and C\'ordoba; they follow from the table on p.~661
of~\cite{BC}, or p.~24 of~\cite{C}. The transition point $x=\nu/2$ in (\ref{1b})
and (\ref{2b}) may be replaced (clearly, with a different $D$) by
$(1+\varepsilon)\nu/2$, or in (\ref{3b}) and (\ref{4b}), $x=1/2$ may be replaced by $(1+\varepsilon)/2$,
where $0<\varepsilon<1$ is chosen
earlier; cf.~\cite{St} for additional comments.

Denoting by $\Phi_\nu(x)$ the function that appears on the right of (\ref{1b})
when $\nu\ge1$ or on the right of (\ref{3b}) when $-1<\nu<1$, we have
\begin{equation}
\label{ec:Jalpha}
|\sqrt{x}J_\nu(x)|\le D \Phi_\nu(x).
\end{equation}
Accordingly, (\ref{2b}) and (\ref{4b}) may be written in the following form
\begin{equation}
\label{ec:cotaJ'}
|\sqrt{x}J'_\nu(x)| \le C \,\frac{\sqrt{|\nu^2-x^2|+x^{4/3}}}{x}\,\Phi_\nu(x)
\end{equation}
when $\nu\ge1$ and
\begin{equation}
\label{ec:cotaJ''}
|\sqrt{x}J'_\nu(x)|\le C \, \frac{x+1}x \,\Phi_\nu(x)
\end{equation}
when $-1<\nu<1$ with the exception for $\nu=0$; then $(x+1)/x$ is
replaced by $x/(x+1)$.

For the purpose of kernel estimates that are proved in Lemma~\ref{Lem:L1}
it is much more convenient to split $(\nu/2,2\nu)$ into the three intervals
$(\nu/2,\nu-\nu^{1/3})$, $(\nu-\nu^{1/3},\nu+\nu^{1/3})$ and $(\nu+\nu^{1/3},
2\nu)$ and, instead of $\Phi_\nu(x)$ in (\ref{ec:Jalpha}) use, if $\nu\ge3$,
say, the (equivalent) bound by
\begin{equation*}
\widetilde{\Phi}_\nu(x) =
\begin{cases}
x^{\nu+1/2}(\Gamma(\nu+1)2^{\nu})^{-1}, & 0<x\le \nu/2,\\
\nu^{1/4}(\nu-x)^{-1/4}, & \nu/2<x\le \nu-\nu^{1/3},\\
\nu^{1/6}, & \nu-\nu^{1/3}<x\le \nu+\nu^{1/3},\\
\nu^{1/4}(x-\nu)^{-1/4}, & \nu+\nu^{1/3}<x\le 2\nu,\\
1,& 2\nu<x.
\end{cases}
\end{equation*}
We have, for certain $a_\nu$, $b_\nu$ and $c_\nu$, the asymptotic
\begin{equation}
\sqrt{t}J_\nu(t)=\sqrt{2/\pi}\bigg(\cos(t+a_\nu)+b_\nu\frac{\sin(t+c_\nu)}{t}+
O(t^{-2})\bigg),\qquad\,t\to\infty.
\label{asy2}
\end{equation}
At several places of the next section, without further refering to it, Stirling's
formula is used:
$$
\lim_{\nu\to\infty}\nu^{\nu+1/2}e^{-\nu}\Gamma(\nu+1)^{-1}=
(2\pi)^{-1/2}.
$$
Given $p$, $1\le p\le\infty$, $p'$ denotes its conjugate, $1/p+1/p'=1$.

%%%%%%%%%%%%%%%%%%%%%%%%%%%%%%%%%%%%%%%%%%%%%%%%%%%%%%
\section{Kernel estimates and bounds for $K_\nu^1$ in different regions}
%%%%%%%%%%%%%%%%%%%%%%%%%%%%%%%%%%%%%%%%%%%%%%%%%%%%%%

In the sequel we write $\overnu=\nu$ if $\nu\ge1$ and
$\overnu=1$ if $-1<\nu\le1$. Also, we use the notation
$$
w_{a,b}(x)=x^a(1+x)^{b-a}.
$$
The lemmas that follow give proper estimates of the kernel $K_\nu^1(x,y)$
in different regions of the quarter plane $(0,\infty)\times(0,\infty)$.
The corresponding propositions furnish weighted $L^p$ bounds of associated
kernel operators restricted to relevant regions with minimal assumptions
required on $a, b, A, B$.
In the propositions the constant $C$ will depend on $\alpha, p, a, b, A, B$
but will not depend on $\nu\ge\alpha$ and~$f$. At several places the
usual interpretation of the $L^p$ norm is needed when $p=\infty$.

%%%%%%%%%%%%%
\begin{Lem}
\label{Lem:L1}
There exists $C>0$ independent of $\nu>-1$ such that
\begin{equation*}
\left|K_\nu^1(x,y)\right|\le C
\begin{cases}
(xy)^{\nu+1/2}/(\Gamma(\nu+2)2^{\nu})^2,& 0<x,y<2\overnu/3,\\
1,& x,y>\overnu/2.
\end{cases}
\end{equation*}
\end{Lem}
%%%%%%%%%%%%%

\begin{proof}
Using (\ref{1b}) and (\ref{3b}) gives, for $0<x,y<2\overnu/3$,
\begin{equation*}
\left|K_\nu^1(x,y)\right| \le C
\frac{(xy)^{\nu+1/2}}{\Gamma(\nu+1)^2 2^{2\nu}} \int_0^1 (1-s^2)
s^{2\nu+1}\,ds.
\end{equation*}
To get the required bound we use
\begin{equation*}
\int_0^1 (1-s^2)s^{2\nu+1}\,ds =
\frac12\frac{\Gamma(\nu+1)\Gamma(2)}{\Gamma(\nu+3)}
\end{equation*}
and then majorize the last expression by $C(1+\nu)^{-2}$.
To prove the estimate for $x,y>\overnu/2$ we note that
by Schwarz's inequality we have $|K_\nu^1(x,y)|^2 \le A_{\nu}(x)A_{\nu}(y)$,
where
$$
A_{\nu}(x) = \int_0^1(1-s^2)(\sqrt{xs}J_{\nu}(xs))^2\,ds.
$$
It is therefore sufficient to check that $A_{\nu}(x)\le C$
for $x>\overnu/2$. In what follows we consider the case $\nu\ge 1$
only (analysing the case $-1<\nu<1$ is much easier). In fact we also assume
that $\nu\ge3$.
We consider four cases to estimate $A_\nu(x)$: $\nu/2<x\le \nu-\nu^{1/3}$,
$\nu-\nu^{1/3}<x\le \nu+\nu^{1/3}$, $\nu+\nu^{1/3}<x\le 2\nu$, and
$2\nu<x$.

\textit{Case 1}: $\nu/2<x\le \nu-\nu^{1/3}$. By (\ref{ec:Jalpha}), (here
and in the sequel we use (\ref{ec:Jalpha}) with $\widetilde{\Phi}_\nu(x)$
in place of $\Phi_\nu(x)$)
\begin{equation*}
|\sqrt{xs}J_\nu(xs)|\le C
\begin{cases}
(xs)^{\nu+1/2}(\Gamma(\nu+1)2^{\nu})^{-1},& 0<s\le \nu/(2x),\\
\nu^{1/4}(\nu-xs)^{-1/4}, & \nu/(2x)<s< 1.
\end{cases}
\end{equation*}
Therefore,
\begin{multline*}
\qquad
A_\nu(x) \le C\bigg(x^{2\nu+1}(\Gamma(\nu+1)2^{\nu})^{-2}
\int_0^{\nu/(2x)}(1-s^2) s^{2\nu+1}\, ds\\
+ \nu^{1/2}\int_{\nu/(2x)}^1(1-s^2) (\nu-xs)^{-1/2}\,ds \bigg).
\qquad
\end{multline*}
The first summand above is bounded by using
\begin{equation*}
\int_0^{\nu/(2x)}(1-s^2) s^{2\nu+1}\, ds\le C
\frac1{2(\nu+1)}\left(\frac{\nu}{2x}\right)^{2(\nu+1)}
\end{equation*}
and Stirling's formula. For the second summand we write
\begin{align*}
\int_{\nu/(2x)}^1(1-s^2) (\nu-xs)^{-1/2}\,ds
&\le \int_{\nu/(2x)}^{\nu/x}(\nu-xs)^{-1/2}\, ds
\end{align*}
and note that the last integral is bounded by $C\nu^{-1/2}$ which gives the
correct bound of the second summand. The claim,
$A_\nu(x)\le C$, now follows.

\textit{Case 2}: $\nu-\nu^{1/3}<x<\nu+\nu^{1/3}$.
From (\ref{ec:Jalpha}), the following estimate holds:
\begin{equation*}
|\sqrt{xs}J_\nu(xs)|\le C
\begin{cases}
(xs)^{\nu+1/2}(\Gamma(\nu+1)2^{\nu})^{-1}, & 0<s\le \nu/(2x),\\
\nu^{1/4}(\nu-xs)^{-1/4}, & \nu/(2x) < s\le (\nu-\nu^{1/3})/x,\\
\nu^{1/6}, & (\nu-\nu^{1/3})/x<s<1.
\end{cases}
\end{equation*}
Therefore,
\begin{multline*}
\quad
A_\nu(x)
\le C\bigg(x^{2\nu+1}(\Gamma(\nu+1)2^{\nu})^{-2}
\int_0^{\nu/(2x)}(1-s^2) s^{2\nu+1}\, ds\\
+ \nu^{1/2}\int_{\nu/(2x)}^{(\nu-\nu^{1/3})/x}
(1-s^2) (\nu-xs)^{-1/2}\,ds +
\nu^{1/3}\int_{(\nu-\nu^{1/3})/x}^1 (1-s^2)\,ds\bigg).
\quad
\end{multline*}
The first summand is that from~\textit{Case~1}. The same is with the second
summand once we enlarge the upper integral limit from $(\nu-\nu^{1/3})/x$
to~$1$.
The integral contained in the third summand is easily seen to be bounded by
$C\nu^{-1/3}$ which gives the correct bound of the summand. The claim,
$A_\nu(x)\le C$, again follows.

\textit{Case 3}: $\nu+\nu^{1/3}<x\le 2\nu$.
As in the previous cases we use (\ref{ec:Jalpha}) obtaining:
\begin{equation*}
|\sqrt{xs}J_\nu(xs)|\le C
\begin{cases}
(xs)^{\nu+1/2}(\Gamma(\nu+1)2^{\nu})^{-1},
& 0<s\le \nu/(2x),\\
\nu^{1/4}(\nu-xs)^{-1/4},
& \nu/(2x) < s \le (\nu-\nu^{1/3})/x,\\
\nu^{1/6},
& (\nu-\nu^{1/3})/x < s \le (\nu+\nu^{1/3})/x,\\
\nu^{1/4}(xs-\nu)^{-1/4},
& (\nu+\nu^{1/3})/x<s<1.
\end{cases}
\end{equation*}
Next, we bound $A_\nu(x)$ by splitting the integration accordingly. The first
two resulting terms are exactly those from~\textit{Case~2}. The same is with
the third one once the upper integral limit is enlarged from
$(\nu+\nu^{1/3})/x$ to~$1$. The fourth resulting term is
\begin{equation*}
\nu^{1/2}\int_{(\nu+\nu^{1/3})/x}^1(1-s^2) (xs-\nu)^{-1/2}\,ds
\end{equation*}
and the required claim follows since
\begin{equation*}
\int_{(\nu+\nu^{1/3})/x}^1(1-s^2)(xs-\nu)^{-1/2}\,ds\le Cx^{-1/2}
\int_{\nu/x}^1\Big(s-\frac{\nu}x\Big)^{-1/2}\,ds
\le C \nu^{-1/2}.
\end{equation*}

\textit{Case 4}: $2\nu<x$. In this situation we have
\begin{equation*}
|\sqrt{xs}J_\nu(xs)| \le C
\begin{cases}
(xs)^{\nu+1/2}(\Gamma(\nu+1)2^{\nu})^{-1},
& 0 < s \le \nu/(2x),\\
\nu^{1/4}(\nu-xs)^{-1/4},
& \nu/(2x) < s \le (\nu-\nu^{1/3})/x,\\
\nu^{1/6},
& (\nu-\nu^{1/3})/x < s \le(\nu+\nu^{1/3})/x,\\
\nu^{1/4}(xs-\nu)^{-1/4},
& (\nu+\nu^{1/3})/x < s \le 2\nu/x,\\
1, & 2\nu/x<s<1.
\end{cases}
\end{equation*}
Again, we bound $A_\nu(x)$ by properly splitting the integration and noting
that the first three resulting terms are those from~\textit{Case~3}. The same is with
the fourth one after replacing the upper integral limit $2\nu/x$ by $1$.
The fifth resulting term is $\int_{2\nu/x}^1(1-s^2)\,ds$ and it is
obviously bounded by a constant. This finishes the proof of~\textit{Case~4}
hence the lemma.
\end{proof}

%%%%%%%%%%%%%
\begin{Prop}
\label{Prop:P1}
Let $\alpha>-1$ and $1\le p\le\infty$. Assuming $a> -1/p-(\alpha+1/2)$
\,$(\ge$ if $p=\infty)$ and $A<1-1/p+(\alpha+1/2)$ \,$(\le$ if $p=1)$
we have
\begin{equation}
\label{ec:xyz} \int_0^{2\overnu/3}
\bigg(w_{a,b}(x)\int_0^{2\overnu/3}
|K_\nu^1(x,y)f(y)|\,dy\bigg)^p\,dx \le C
\int_0^{2\overnu/3}\left|f(y)w_{A,B}(y)\right|^p\,dy,
\end{equation}
for $\nu\ge\alpha$, while, assuming in addition $b\le B$, for $\nu\ge \alpha$
we have
\begin{equation}
\label{ec:diagonal} \int_{\overnu/2}^{\infty}
\bigg(w_{a,b}(x) \int_{\overnu/2}^{\infty}\chi_{D}(x,y)
|K_\nu^1(x,y)f(y)|\,dy\bigg)^p\,dx \le
C\int_{\overnu/2}^{\infty}\left|f(y)w_{A,B}(y)\right|^p\,dy,
\end{equation}
where $D=\{(x,y): |x-y|\le 1\}$.
\end{Prop}
%%%%%%%%%%%%%

\begin{proof}
Proving (\ref{ec:xyz}) use the bound on $|K_\nu^1(x,y)|$ from
Lemma~\ref{Lem:L1} and consider $p=1$ first. Then split the integration
regions $[0,2\overnu/3]$ in the double integral using $2/3$ as
breakpoints and change
the order of integration to end up with four integrals to be estimated by
a constant multiple of either
$\int_0^{2/3}|f(y)|y^A\,dy$ or $\int_{2/3}^{2\overnu/3}|f(y)|y^B\,dy$.
If the $y$-region of integration is $[0,2/3]$, then merely replace $y^{\nu+1/2}$
by $y^A$ (using the assumption $A\le \alpha+1/2$) and note that the integrals
\begin{equation*}
\int_0^{2/3}x^{a+\nu+1/2}\,dx,
\qquad \int_{2/3}^{2\overnu/3} x^{b+\nu+1/2}\,dx,
\end{equation*}
when multiplied by $(\Gamma(\nu+2)2^{\nu})^{-2}$ give
bounded functions of $\nu\ge\alpha$. If the $y$-region of integration is $[2/3,
2\overnu/3]$ then, for large values of $\nu$ say, replace $y^{\nu+1/2}$
by $(2\nu/3)^{\nu+1/2-B}y^B$ and, in the case of the $x$-integration over
$[2/3,2\nu/3]$, note that
\begin{equation*}
\frac1{(\Gamma(\nu+2)2^{\nu})^2}
\,\Big(\frac23\nu\Big)^{\nu+1/2-B}\int_{2/3}^{2\nu/3} x^{b+\nu+1/2}\,dx
\end{equation*}
is a bounded function of large $\nu$ (note that the possible growth of
$x^{b-B-2}$ is compensated by the decay of $(e/3)^\nu$).
The remaining case of the $x$-integration over $[0,2/3]$ is similar (and easier).
In the case $1<p<\infty$ apply the analogous argument using H\"older's
inequality in appropriate places (cf.\ also the proof of Lemma~2.1
in~\cite{St}). Consider now $p=\infty$.
It is then sufficient to show that, for $\nu\ge\alpha$, the quantities
\begin{equation*}
\frac{1}{(\Gamma(\nu+2)2^{\nu})^2}\sup_{0<x<2/3}\bigg[x^{a+\nu+1/2}
\int_0^{2\overnu/3}y^{\nu+1/2}|f(y)|\,dy\bigg]
\end{equation*}
and
\begin{equation*}
\frac{1}{(\Gamma(\nu+2)2^{\nu})^2}\sup_{2/3<x<2\overnu/3}\bigg[x^
{b+\nu+1/2}\int_0^{2\overnu/3}y^{\nu+1/2}|f(y)|\,dy\bigg]
\end{equation*}
are bounded by a constant (independent of $\nu\ge\alpha$) multiplied by the sum
\begin{equation*}
\sup_{0<y<2/3}\,y^A|f(y)| + \sup_{2/3<y<2\overnu/3}y^B|f(y)|.
\end{equation*}
This is easily done by using the assumptions $a\ge-(\alpha+1/2)$ and
$A<1+(\alpha+1/2)$.

Proving (\ref{ec:diagonal}) use Lemma~\ref{Lem:L1} and note that only
$\nu\ge1$ may be considered. Let $p=1$.
It is sufficient to show that
\begin{equation*}
\int_{\nu/2}^{\infty}x^{b}\int_{\nu/2}^{\infty}\chi_{[0,1]}(|x-y|)|f(y)|\,dy
\,dx\le C \int_{\nu/2}^{\infty}\left|f(y)\right|y^B\,dy
\end{equation*}
with $C$ independent of $\nu\ge1$.
The inequality immediately follows by changing the order of integration,
noting that for $y\ge\nu/2$ and $\nu\ge1$
\begin{equation*}
\int_{\nu/2}^{\infty}x^{b}\chi_{[0,1]}(|x-y|)\,dx
\le Cy^b,
\end{equation*}
and using $b\le B$. Consider now the case $1<p\le\infty$. By using
$b\le B$ it is easily seen that (\ref{ec:diagonal}) is implied by
\begin{equation*}
\int_{\nu/2}^{\infty}\left(\int_{x-1}^{x+1}|
\chi_{[\nu/2,\infty]}(y)g(y)|\,dy\right)^p\,dx
\le C \int_{\nu/2}^{\infty}\left|g(y)\right|^p\,dy,
\end{equation*}
where $\nu\ge1$.
This, in order, follows from the well-known standard norm inequality for the
maximal function since, for $x\ge\nu/2$ and $\nu\ge1$,
\begin{equation*}
\int_{x-1}^{x+1}|\chi_{[\nu/2,\infty]}(y)g(y)|\,dy
\le C M(\chi_{[\nu/2,\infty]}g)(x).
\end{equation*}
\end{proof}

%%%%%%%%%%%%%
\begin{Lem}
\label{Lem:L2} Let $\alpha>-1$. There exists a constant $C$
independent of $\nu\ge\alpha$ such that for $|x-y|\ge1$ and
$x,y\ge 5\overnu/4$,
\begin{equation}
\label{ec:cota2nu}
|K^1_\nu(x,y)|\le C |x-y|^{-2}.
\end{equation}
\end{Lem}
%%%%%%%%%%%%%

\begin{proof}
Using (\ref{Cam}) gives
\begin{equation}
\label{ec:jk}
|K^1_\nu(x,y)|\le 2 |u|^{-2}\bigg(|2\nu^2-v||F_1|+\frac{2v}{|u|}
(x|F_2|+y|F_3|)+2xy|F_4|\bigg).
\end{equation}
Since $|F_i|\le C$ on $x,y\ge\overnu$, see (\ref{1b})--(\ref{4b}),
the conclusion then follows by noting that each of the terms,
$|2\nu^2-(x^2+y^2)|$, $v$, $xy$, when
divided by $(x+y)^{2}$ is bounded on the indicated range of $x$ and~$y$.
\end{proof}

%%%%%%%%%%%%%
\begin{Prop}
\label{Prop:P2}
Let $\alpha>-1$, $1\le p\le\infty$, $b\le B$, $b<2-1/p$ $(\le$ if $p=\infty)$,
and $B>-1-1/p$ $(\ge$ if $p=1)$. Then,
for $\nu\ge \alpha$,
\begin{equation}
\label{ec:antidiagonal} \int_{5\overnu/4}^{\infty}
\bigg(w_{a,b}(x) \int_{5\overnu/4}^{\infty}
\chi_{D^c}(x,y)|K_\nu^{1}(x,y)f(y)|\,dy \bigg)^p\,dx \le
C\int_{5\overnu/4}^{\infty}
\left|f(y)w_{A,B}(y)\right|^p\,dy,
\end{equation}
where $D^c=\{(x,y): |x-y|\ge 1\}$.
\end{Prop}
%%%%%%%%%%%%%

\begin{proof}
Consider first $p=1$. Then (\ref{ec:antidiagonal}) reduces to showing that
\begin{equation*}
\int_{5\overnu/4}^{\infty}
\int_{5\overnu/4}^{\infty}\chi_{[1,\infty)}
(|x-y|)x^by^{-B}\frac{|g(y)|}{|x-y|^2}\,dy\,dx
\le C\int_{2\overnu}^{\infty}|g(y)|\,dy
\end{equation*}
and this is a consequence of (this time $b<1$ is used)
\begin{equation*}
\sup_{\nu\ge\alpha}\sup_{y\ge5\overnu/4}\bigg[y^{-B}
\int_{5\overnu/4}^{\infty}x^b\chi_{[1,\infty)}
(|x-y|)\frac{1}{|x-y|^2}\,dx\bigg] < \infty.
\end{equation*}
In the case $p=\infty$, (\ref{ec:antidiagonal}) reduces to
showing that
\begin{equation*}
\sup_{x\ge5\overnu/4}
\bigg[x^b\int_{5\overnu/4}^{\infty}\chi_{D^c}(x,y)\frac{|g(y)|}
{|x-y|^2}\,y^{-B}\,dy\bigg] \le C\,\sup_{y\ge5\overnu/4}|g(y)|
\end{equation*}
and this easily follows by noting that ($B>-1$ is used)
\begin{equation*}
\sup_{\overnu\ge\alpha}\sup_{x\ge5\overnu/4}
\bigg[x^b\int_{2\overnu}^{\infty}
\chi_{[1,\infty)}(|x-y|)\frac{1}{y^B|x-y|^2}\,dy\bigg]<\infty.
\end{equation*}
Consider now the case $1<p<\infty$. We decompose the $y$-range in
$[5\overnu/4, x/2]$, $[x/2,2x]$ and $[2x,\infty)$. Then,
\begin{multline*}
\qquad
\int_{5\overnu/4}^{\infty} \bigg(w_{a,b}(x)
\int_{5\overnu/4}^{x/2}
\chi_{D^c}(x,y) |K_\nu^1(x,y)f(y)|\,dy\bigg)^p\,dx\\
\le C\int_{5\overnu/4}^\infty \bigg(x^{b-2}
\int_{5\overnu/4}^{x}|f(y)|\,dy\bigg)^p\,dx,
\qquad
\end{multline*}
and an application of H\"older's inequality shows that the last expression is
bounded by the right side of~(\ref{ec:antidiagonal}). Similarly,
\begin{equation*}
\int_{5\overnu/4}^{\infty} \bigg(w_{a,b}(x)\int_{2x}^{\infty}
\chi_{D^c}(x,y) |K_\nu^{1}(x,y)f(y)|\,dy\bigg)^p\,dx
\le C\int_{5\overnu/4}^\infty \bigg(x^{b}\int_{x}^{\infty}
\frac{|f(y)|}{y^2}\, dy\bigg)^p\,dx
\end{equation*}
and again H\"older's inequality does the job. Proving
\begin{equation*}
\int_{5\overnu/4}^{\infty} \bigg(w_{a,b}(x)\int_{x/2}^{2x}
\chi_{D^c}(x,y) |K_\nu^{1}(x,y)f(y)|\,dy \bigg)^p\,dx
\le C \int_{5\overnu/4}^{\infty}|w_{A,B}(x)f(x)|^p\,dx
\end{equation*}
reduces to showing that
\begin{equation*}
\int_{5\overnu/4}^{\infty}x^{p(b-B)}\bigg(\int_{x/2}^{2x}
\chi_{[1,\infty]}(|x-y|)\frac{|g(y)|}{|x-y|^2}\,dy\bigg)^p\,dx
\le C \int_{5\overnu/4}^{\infty}|g(x)|^p\,dx.
\end{equation*}
By the assumption $b\le B$, we obtain
\begin{align*}
&\phantom{\null\le\null}
\int_{5\overnu/4}^{\infty}x^{p(b-B)}\bigg(\int_{x/2}^{2x}
\chi_{[1,\infty]}(|x-y|)\frac{|g(y)|}{|x-y|^2}\,dy\bigg)^p\,dx\\
&\le C\int_{5\overnu/4}^\infty
\left( \sum_{k=1}^\infty\int_{|x-y|\sim 2^k} \frac{|g(y)|}{|x-y|^2}
\chi_{[5\overnu/4,\infty]}(y)\, dy\right)^p\, dx\\
&\le C\int_{5\overnu/4}^\infty
\left(\sum_{k=1}^\infty 2^{-2k}\int_{|x-y|\sim2^k} |g(y)|
\chi_{[5\overnu/4,\infty]}(y)\, dy\right)^p\,dx\\
&\le C\int_{5\overnu/4}^\infty
M(\chi_{[5\overnu/4,\infty)}g)(x)^p\,dx.
\end{align*}
Using the fact that $M$ is a bounded operator on $L^p$ for $1<p<\infty$,
we conclude the proof of Proposition~\ref{Prop:P2}.
\end{proof}

%%%%%%%%%%%%%
\begin{Lem}
\label{Lem:L3}
There exists a constant $C$ independent of $\nu\ge1$ such that for
$|x-y|\ge 1$ and $\nu/2\le x,y\le 3\nu/2$,
\begin{equation}
\label{ec:cotanumedios2nu}
|K^1_\nu(x,y)|\le C |x-y|^{-2}
\left(\left(\frac{|\nu-x|+\nu^{1/3}}{|\nu-y|+\nu^{1/3}}\right)^{1/4}
+ \left(\frac{|\nu-y|+\nu^{1/3}}{|\nu-x|+\nu^{1/3}}\right)^{1/4}\right).
\end{equation}
\end{Lem}
\begin{proof}
%%%%%%%%%%%%%

Applying (\ref{ec:Jalpha}) and (\ref{ec:cotaJ'}) to (\ref{ec:jk}) leads to
\begin{multline}
\label{ec:basic}
|K_\nu^1(x,y)|\le C \Phi_\nu(x)\Phi_\nu(y)|u|^{-2}\,
\bigg(|2\nu^2-v|+\frac{v}{|u|}(|\nu^2-x^2|+x^{4/3})^{1/2}\\
+ \frac{v}{|u|}(|\nu^2-y^2|+y^{4/3})^{1/2}
+ (|\nu^2-x^2|+x^{4/3})^{1/2}(|\nu^2-y^2|+y^{4/3})^{1/2}\bigg).
\end{multline}
Then the estimates for
$\Phi_\nu$ on the range $(\nu/2,3\nu/2)$ give
\begin{multline*}
|K_\nu^1(x,y)|\le C|x-y|^{-2}(|\nu-x|+\nu^{1/3})^{-1/4}(|\nu-y|+\nu^{1/3})
^{-1/4}\nu^{-3/2}\\ \times
\bigg(\nu\big[|\nu-x|+|\nu-y|\big]+\nu^{3/2}\big[(|\nu-x|+\nu^{1/3})^{1/2}+
(|\nu-y|+\nu^{1/3})^{1/2}\big]\\
+\nu\big(|\nu-x|+\nu^{1/3}\big)^{1/2}\big(|\nu-y|+\nu^{1/3}\big)^{1/2}\bigg).
\end{multline*}
A careful analysis then shows that each of the three resulting summands is
bounded by the right side of~(\ref{ec:cotanumedios2nu}).
\end{proof}

%%%%%%%%%%%%%
\begin{Prop}
\label{Prop:P3}
Let $\alpha\ge1$, $1\le p\le\infty$, and $b\le B$. Then, for $\nu\ge \alpha$,
\begin{equation}
\label{ec:d2}
\int_{\nu/2}^{3\nu/2} \bigg(w_{a,b}(x)
\int_{\nu/2}^{3\nu/2}\chi_{D^c}(x,y)|K_\nu^{1}(x,y)f(y)|\,dy\bigg)^p dx
\le C\int_{\nu/2}^{3\nu/2}\left|f(y)w_{A,B}(y)\right|^p dy.
\end{equation}
\end{Prop}
%%%%%%%%%%%%%

\begin{proof}
By using (\ref{ec:cotanumedios2nu}) and the assumption $b\le B$,
(\ref{ec:d2}) will follow from
\begin{multline}
\label{ec:abc}
\quad
\bigg(\int_{\nu/2}^{3\nu/2}\bigg(\int_{\nu/2}^{3\nu/2}\chi_{D^c}(x,y)\frac{|f(y)|}
{|x-y|^2}\bigg(\frac{|\nu-x|+\nu^{1/3}}{|\nu-y|+\nu^{1/3}}\bigg)^{1/4}\,dy
\bigg)^p\,dx\bigg)^{1/p} \\
\le C\bigg(\int_{\nu/2}^{3\nu/2}|f(y)|^p\,dy\bigg)^{1/p}
\qquad
\end{multline}
and an analogous inequality with the second term from the right side of
(\ref{ec:cotanumedios2nu}) involved. Since in the latter case the argument is
completely similar to what we present below we do not repeat it.
The change of variables $X=\nu^{-1/3}(x-\nu)$,
$Y=\nu^{-1/3}(y-\nu)$, shows that the left side of (\ref{ec:abc}) equals
$\nu^{-1/3+1/(3p)}$ multiplied by
\begin{equation*}
\bigg(\int_{-\nu^{2/3}/2}^{\nu^{2/3}/2}\bigg(\int_{-\nu^{2/3}/2}
^{\nu^{2/3}/2}\chi_{\{|X-Y|>\nu^{-1/3}\}}(X,Y)\frac{|F(Y)|}{|X-Y|^2}
\bigg(\frac{1+|X|}{1+|Y|}\bigg)^{1/4}\,dY\bigg)^p\,dX\bigg)^{1/p},
\end{equation*}
with $F(Y)=f(Y\nu^{1/3}+\nu)$. By symmetry, it is now enough to majorize the
quantity
\begin{equation*}
\bigg(\int_{0}^{\nu^{2/3}/2}\bigg(\int_{-\nu^{2/3}/2}
^{\nu^{2/3}/2}\chi_{\{|X-Y|>\nu^{-1/3}\}}(X,Y)\frac{|F(Y)|}{|X-Y|^2}
\bigg(\frac{1+|X|}{1+|Y|}\bigg)^{1/4}\,dY\bigg)^p\,dX\bigg)^{1/p}
\end{equation*}
by $\nu^{1/3-1/(3p)}$ times the right side of (\ref{ec:abc}).
We decompose the $Y$-range into five intervals $[-\nu^{2/3}/2,-X]$,
$[-X,0]$, $[0,X/2]$, $[X/2,2X]$ and $[2X,\nu^{2/3}/2]$ (the extreme intervals
may be void, the second one may reduce to $[-\nu^{2/3}/2,0]$) and denote
the resulting double integrals by $I_i$, $i=1,\dots,5$. We will focus on
proving that
\begin{equation}
\label{ec:bbb}
I_i\le C\nu^{1/3-1/(3p)}\bigg(\int_{\nu/2}^{3\nu/2}|f(y)|^p\,dy
\bigg)^{1/p},\qquad i=1,\dots,5.
\end{equation}
Let $i=1$. If $Y\in [-\nu^{2/3}/2, -X]$ then $|X-Y|^{-2}\le C|Y|^{-2}$,
the change of variable $Y=-T$ and H\"older's inequality produce the following
bound of~$I_1$:
\begin{align*}
&\phantom{\null\le\null}
\bigg(\int_{0}^{\nu^{2/3}/2} \bigg(\int_{-\nu^{2/3}/2}^{-X}
\chi_{\{|X-Y|>\nu^{-1/3}\}}(X,Y)\frac{|F(Y)|}{|X-Y|^2}
\bigg(\frac{1+|X|}{1+|Y|}\bigg)^{1/4}\,dY\bigg)^p\,dX\bigg)^{1/p}\\
&\le C\bigg(\int_{0}^{\nu^{2/3}/2}\bigg(\int_{X}^{\nu^{2/3}/2}
\chi_{\{X+T>\nu^{-1/3}\}}(X,T)\frac{|F(-T)|}{T^{2}}
\bigg(\frac{1+X}{1+T}\bigg)^{1/4}\,dT\bigg)^p\,dX\bigg)^{1/p}\\
&\le C\bigg(\int_{0}^{\nu^{2/3}/2}\bigg(\int_{\max\{X,\nu^{-1/3}-X\}}
^{\nu^{2/3}/2} \frac{|F(-T)|}{T^{2}}\,dT\bigg)^p\,dX\bigg)^{1/p}\\
&\le C\nu^{-1/(3p)}\bigg(\int_{0}^{\nu^{2/3}/2} \bigg(\int_{\max\{X,2\nu^
{-1/3}-X\}}^{\nu^{2/3}/2}T^{-2p'}\,dT\bigg)^{p/p'}dX\bigg)^{1/p} \\
&\hspace{0.6\textwidth} \times
\bigg(\int_{\nu/2}^{3\nu/2}|f(y)|^p\,dy\bigg)^{1/p}.
\end{align*}
This proves (\ref{ec:bbb}) for $i=1$ since
\begin{equation*}
\bigg(\int_{0}^{\nu^{2/3}/2}\bigg(\int_{\max\{X,\nu^{-1/3}-X\}}^
{\nu^{2/3}/2}T^{-2p'}\,dT\bigg)^{p/p'}dX\bigg)^{1/p}\le C\nu^{1/3}.
\end{equation*}
Let $i=2$. If $Y\in [-X,0]$ then $|X-Y|^{-2}\le C|X|^{-2}$, the change of
variable $Y=-T$ and H\"older's inequality similarly show that $I_2$ is bounded
by $\nu^{-1/(3p)}$ multiplied by the product of
\begin{equation}
\label{ec:w1}
\bigg(\int_{0}^{\nu^{2/3}/2}X^{-2p}\bigg(\int_{\max\{0,
\nu^{-1/3}-X\}}^{X}\bigg(\frac{1+X}{1+T}\bigg)^{p'/4}\,dT\bigg)^{p/p'}
dX\bigg)^{1/p}
\end{equation}
and the right side of~(\ref{ec:abc}). This completes the proof of
(\ref{ec:bbb}) for~$i=2$ once we note that (\ref{ec:w1}) is~$O(\nu^{1/3})$.
Consider the case~$i=3$. Then $I_3$ equals
\begin{equation*}
\bigg(\int_{\nu^{-1/3}}^{\nu^{2/3}/2}\bigg(\int_{0}^{X/2}
\chi_{\{|X-Y|>\nu^{-1/3}\}}(X,Y)\frac{|F(Y)|}{|X-Y|^2}
\bigg(\frac{1+X}{1+Y}\bigg)^{1/4}\,dY\bigg)^p\,dX\bigg)^{1/p}
\end{equation*}
and we further split the $X$-range onto $[\nu^{-1/3},1]$ and $[1,\nu^{2/3}/2]$
denoting the resulting double integrals by $I_{31}$ and $I_{32}$. Then
\begin{align*}
I_{31}
&\le C\bigg(\int_{\nu^{-1/3}}^{1}\bigg(\int_{0}^{\min\{X/2,X-\nu^{-1/3}\}}
\frac{|F(Y)|}{X^{2}}\,dY\bigg)^p\,dX\bigg)^{1/p}\\
&\le C\bigg[\bigg(\int_{\nu^{-1/3}}^{2\nu^{-1/3}}\bigg(\int_{0}^{X-\nu^{-1/3}}
\frac{|F(Y)|}{X^{2}}\,dY\bigg)^p\,dX\bigg)^{1/p}\\
&\qquad \qquad \qquad
+\bigg(\int_{2\nu^{-1/3}}^{1}\bigg(\int_{0}^{X/2}\frac{|F(Y)|}{X^{2}}
\,dY\bigg)^p\,dX\bigg)^{1/p}\bigg]\\
&\le C\nu^{-1/(3p)}\bigg[\bigg(\int_{\nu^{-1/3}}^{2\nu^{-1/3}}X^{-2p}
(X-\nu^{-1/3})^{p/p'}\,dX\bigg)^{1/p}\\
&\qquad \qquad \qquad
+\bigg(\int_{2\nu^{-1/3}}^{1}X^{p(1/p'-2)}\,dX\bigg)^{1/p}\bigg]
\bigg(\int_{\nu/2}^{3\nu/2}|f(y)|^p\,dy\bigg)^{1/p}
\end{align*}
where, in the last line, H\"older's inequality was used. Now, the bound
\begin{equation*}
\bigg(\int_{\nu^{-1/3}}^{2\nu^{-1/3}}X^{-2p}(X-\nu^{-1/3})^{p/p'}dX\bigg)^{1/p} +
\bigg(\int_{2\nu^{-1/3}}^{1}X^{p(1/p'-2)}\,dX\bigg)^{1/p}\le C \nu^{1/3}
\end{equation*}
shows (\ref{ec:bbb}) for $I_{31}$. Similarly, using H\"older's inequality
again, we majorize $I_{32}$ by $\nu^{-1/(3p)}$ times the product of
\begin{equation}
\label{ec:ddd}
\bigg(\int_{1}^{\nu^{2/3}/2}X^{-7p/4}
\bigg(\int_{0}^{X}(1+Y)^{-p'/4}\,dY\bigg)^{p/p'}dX\bigg)^{1/p}
\end{equation}
and the right side of (\ref{ec:abc}).
Thus (\ref{ec:bbb}) follows for $I_{32}$, hence also for $I_3$, since
(\ref{ec:ddd}) is $O(\nu^{1/3})$.
Consider the case~$i=4$. Then $I_4$ equals
\begin{equation*}
\bigg(\int_{0}^{\nu^{2/3}/2}\bigg(\int_{X/2}^{2X}
\chi_{\{|X-Y|>\nu^{-1/3}\}}(X,Y)\frac{|F(Y)|}{|X-Y|^2}
\bigg(\frac{1+X}{1+Y}\bigg)^{1/4}\,dY\bigg)^p\,dX\bigg)^{1/p}
\end{equation*}
and we further bound it by
\begin{align*}
&\phantom{\null\le\null}
C\bigg(\int_{0}^{\nu^{2/3}/2}\bigg(\int_{X/2}^{2X}\chi_{\{|X-Y|>\nu^
{-1/3}\}}(X,Y)\frac{|F(Y)|}{|X-Y|^2}\,dY\bigg)^p\,dX\bigg)^{1/p}\\
&\le C\bigg(\int_{0}^{\nu^{2/3}/2}\bigg( \sum_{k=k_1}^{\infty}\int_{|X-Y|
\sim 2^k}\frac{|\chi_{[0,\nu^{2/3}/2]}(Y)F(Y)|}{|X-Y|^2}\,dY\bigg)^p\,dX\bigg)^{1/p}\\
&\le C \bigg(\int_{0}^{\nu^{2/3}/2}\bigg(\sum_{k=k_1}^{\infty}2^{-2k}\int_{|X-Y|
\sim 2^k} |\chi_{[0,\nu^{2/3}/2]}(Y)F(Y)|\,dY\bigg)^p\,dX\bigg)^{1/p}\\
&\le C \bigg(\int_{0}^{\nu^{2/3}/2}\bigg(\sum_{k=k_1}^{\infty}2^{-k}
M(\chi_{[0,\nu^{2/3}/2]}F)(X)\bigg)^p \,dX\bigg)^{1/p}\\
&\le C\nu^{1/3}\bigg(\int_{0}^{\nu^{2/3}/2}M(\chi_{[0,\nu^{2/3}/2]}F)
(X)^p\,dX\bigg)^{1/p}
\end{align*}
with $k_1 = [-(\log_2\nu)/3]-1$. Using the fact that $M$ is a bounded
operator on $L^p$ we finish the proof of (\ref{ec:bbb}) for $i=4$.
Consider the case~$i=5$. Then $I_5$ equals
\begin{equation*}
\bigg( \int_{0}^{\nu^{2/3}/2}\bigg(\int_{2X}^{\nu^{2/3}/2}
\chi_{\{|X-Y|>\nu^{-1/3}\}}(X,Y)\frac{|F(Y)|}{|X-Y|^2}
\bigg(\frac{1+X}{1+Y}\bigg)^{1/4}\,dY\bigg)^p \,dX\bigg)^{1/p}
\end{equation*}
and is easily seen to be bounded by a constant multiplied by the sum of
\begin{equation*}
I_{51}=\bigg(\int_{0}^{1}\bigg(\int_{2X}^{\nu^{2/3}/2}
\chi_{\{|X-Y|>\nu^{-1/3}\}}(X,Y)\frac{|F(Y)|}{Y^{2}}\,dY\bigg)^p\,dX\bigg)^{1/p}
\end{equation*}
and
\begin{equation*}
I_{52}=\bigg(\int_{1}^{\nu^{2/3}/2}\bigg(\int_{2X}^{\nu^{2/3}/2}
\chi_{\{|X-Y|>\nu^{-1/3}\}}(X,Y)\frac{|F(Y)|}{Y^{2}}\,dY\bigg)^p\,dX\bigg)^{1/p}.
\end{equation*}
Using H\"older's inequality and the estimate
\begin{equation*}
\bigg(\int_{0}^{1}\bigg(\int_{\max\{X+\nu^{-1/3},2X\}}^{\nu^{2/3}/2}Y^{-2}
\,dY\bigg)^{p/p'}dX\bigg)^{1/p}\le C \nu^{1/3}
\end{equation*}
(decompose the $X$-range into $[0,\nu^{-1/3}]$ and $[\nu^{-1/3},1]$) shows
that
\begin{align*}
I_{51}&\le C \bigg(\int_{0}^{1}\bigg(\int_{\max\{X+\nu^{-1/3},2X\}}
^{\nu^{2/3}/2}\frac{|F(Y)|}{Y^{2}}\,dY\bigg)^p\,dX\bigg)^{1/p}\\
&\le C\nu^{1/3-1/(3p)}\bigg(\int_{\nu/2}^{3\nu/2}|f(y)|^p\,dy\bigg)^{1/p}.
\end{align*}
Similarly, the estimate
\begin{equation*}
\bigg(\int_{1}^{\nu^{2/3}/2}\bigg(\int_{2X}^{\nu^{2/3}/2}
Y^{-2p'}\,dY\bigg)^{p/p'}dX\bigg)^{1/p}\le C\nu^{1/3}
\end{equation*}
and H\"older's inequality produce
\begin{align*}
I_{52}&\le C\bigg(\int_{1}^{\nu^{2/3}/2}
\bigg(\int_{2X}^{3\nu^{2/3}}Y^{-2p'}\,dY\bigg)^{p/p'}dX\bigg)^{1/p}
\bigg(\int_{-\nu^{2/3}/2}^{\nu^{2/3}/2}|F(Y)|^p\,dY\bigg)^{1/p}\\
&\le C\nu^{1/3-1/(3p)}\bigg(\int_{\nu/2}^{3\nu/2}|f(y)|^p\,dy\bigg)^{1/p}.
\end{align*}
This finishes the proof of (\ref{ec:bbb}) for $i=5$ hence the proposition.
\end{proof}

%%%%%%%%%%%%%
\begin{Lem}
\label{Lem:L4}
Let
\begin{center}
\begin{tabular}{rr}
$R_1=[2\overnu/3,\infty)\times[0,\overnu/2]$, &
$\quad R_2=
[3\overnu/2,\infty)\times[\overnu/2,5\overnu/4]$,\\
$R_3=[0,\overnu/2]\times[2\overnu/3,\infty)$, &
$\quad R_4=
[\overnu/2,5\overnu/4]\times[3\overnu/2,\infty)$.\\
\end{tabular}
\end{center}
Then there exists a constant $C$ independent of $\nu>-1$ such that
\begin{equation*}
|K_\nu^1(x,y)\chi_{R_i}(x,y)|\le C \,\frac{\Phi_\nu(x)\Phi_\nu(y)}{|x-y|^2},
\qquad i=1,2,3,4.
\end{equation*}
\end{Lem}
%%%%%%%%%%%%%

\begin{proof}
We discuss the case $\nu\ge1$ only (if $\nu<1$ our argument also applies but
the right side of (\ref{ec:basic}) must be suitably modified).
It is sufficient to check that the expression in parentheses on the right side
of (\ref{ec:basic}) divided by $(x+y)^2$ is bounded by a constant.
In $R_1$, $x\le x+y \le 7x/4$ and $x/4\le x-y\le x$, therefore $v/|u|$ is
bounded. Also each of the terms $(|\nu^2-x^2|+x^{4/3})^{1/2}$,
$(|\nu^2-y^2|+y^{4/3})^{1/2}$, when divided by $x+y$ is bounded. Finally,
$|2\nu^2-(x^2+y^2)|$ when divided by $(x+y)^2$ is bounded. Similar analysis
holds for $i=2,3,4$.
\end{proof}

%%%%%%%%%%%%%
\begin{Prop}
\label{Prop:P4}
Let $\alpha>-1$, $1\le p\le\infty$ and $a,A,b,B$ satisfy the assumptions
$(\ref{t1})$ and $(\ref{t2})$ of Theorem~$\ref{Theorem:main}$. Assume also that
\begin{equation*}
b \le
\begin{cases}
B+1+1/2-2/(3p),& \text{ if } 1\le p<4/3,\\
B+1,& \text{ if } 4/3\le p\le 4,
\,\,\,(< \text{if } p=4/3 \text{ or } p=4),\\
B+1-1/6+2/(3p),& \text{ if } 4<p\le\infty.
\end{cases}
\end{equation*}
Then, for $\nu\ge \alpha$,
\begin{equation}
\label{dod} \int_{0}^{\infty}\bigg(w_{a,b}(x)\int_{0}^{\infty}
\chi_{R_i}(x,y)|K_\nu^{1}(x,y)f(y)|\,dy\bigg)^p\,dx \le
C\int_{0}^{\infty}|f(y)w_{A,B}(y)|^p\,dy,
\end{equation}
where $R_i$, $i=1,\ldots,4$, are as in Lemma~$\ref{Lem:L4}$.
\end{Prop}
%%%%%%%%%%%%%

\begin{proof}
Denote the left side of the inequality to be proved raised to the $1/p$\,th
power (with the usual interpretation when $p=\infty$) by $F_i=F_i(a,b,p,
\nu;f)$. Consider first the case~$i=1$. Using the bound
from Lemma~\ref{Lem:L4} and H\"older's inequality gives
\begin{align*}
F_1 &= \bigg(\int_{0}^{\infty} \left(w_{a,b}(x)\int_{0}^{\infty}
\chi_{R_1}(x,y) |K_{\nu}^{1}(x,y)f(y)|\,dy\right)^p\,dx\bigg)^{1/p}\\
&\le C \bigg(\int_{2\overnu/3}^{\infty}\bigg(x^b \int_{0}^{\overnu/2}
|f(y)|\frac{\Phi_{\nu}(x)
\Phi_{\nu}(y)}{|x-y|^2}\,dy\bigg)^p\,dx\bigg)^{1/p}\\
&\le C \bigg(\int_{2\overnu/3}^{\infty}\bigg(x^{b-2}
\Phi_{\nu}(x)\int_{0}^{\overnu/2} |f(y)|
\Phi_{\nu}(y)\,dy\bigg)^p\,dx\bigg)^{1/p}\\
&\le C \|\chi_{(2\overnu/3,\infty)}(x)x^{b-2}\Phi_{\nu}(x)\|_{p}
\cdot \|\chi_{(0,\overnu/2)}(y)w_{A,B}(y)^{-1}
\Phi_{\nu}(y)\|_{p'} \cdot \|f(y)w_{A,B}(y)\|_p.
\end{align*}
The product of the first two terms following the last $C$ above is a continuous function
of $\nu\in[\alpha,\infty)$. It is therefore sufficient to consider large
$\nu$'s; analogous remark applies when analysing $F_i$, $i=2,3,4$. Thus, from now on
we assume $\nu$ to be large and write $\nu$ in place of $\overnu$.
Then we have
\begin{align}
\label{ec:estimacion0}
&\phantom{\null\le\null}
\|\chi_{(0,\nu/2)}(y)w_{A,B}(y)^{-1}\Phi_\nu(y)\|_{p'}\\
&\le C\max\{\|\chi_{(0,1)}(y)y^{-A}\Phi_\nu(y)\|_{p'},
\|\chi_{(1,\nu/2)}(y)y^{-B}\Phi_\nu(y)\|_{p'}\}\nonumber\\
&\le C 2^{-\nu}\Gamma(\nu+1)^{-1}(\nu/2)^{-B+\nu+1/2}\nonumber,
\end{align}
where, to assure the convergence of the integral on $[0,1]$ the condition
$A<1-1/p+(\alpha+1/2)$ was used ($A\le\alpha+1/2$ when $p=1$). Moreover,
\begin{equation}
\label{ec:normapesob}
\|\chi_{(2\nu/3,\infty)}(x)x^{b-2}\Phi_\nu(x)\|_{p} \le C\nu^{b-2+1/p}
\begin{cases}
1, & \text{ if } p<4,\\
(\log \nu)^{1/4}, & \text{ if } p=4,\\
\nu^{1/6-2/(3p)}, & \text{ if } p>4.\\
\end{cases}
\end{equation}
This is obtained by splitting the integration onto $(2\nu/3,2\nu)$ and
$(2\nu,\infty)$, using the assumption $b<2-1/p$ ($b\le2$ when $p=\infty)$
and then by inspection.
We conclude the analysis by noting that the product of majorants in
(\ref{ec:estimacion0}) and (\ref{ec:normapesob}) produces a bounded function
of $\nu\ge1$.

In the case $i=2$, to get the required bound we proceed as in the case just discussed
majorizing $F_2$ by a constant times
\begin{equation*}
\|\chi_{(3\nu/2,\infty)}(x)x^{b-2}\|_{p}\cdot
\|\chi_{(\nu/2,5\nu/4)}(y)w_{A,B}(y)^{-1}\Phi_\nu(y)\|_{p'}\cdot
\|f(y)w_{A,B}(y)\|_p.
\end{equation*}
Again, we claim that the product of the first two terms
above is a bounded function of $\nu\ge1$. This is because
$\|\chi_{(3\nu/2,\infty)}(x)x^{b-2}\|_{p}=O(\nu^{b-2+1/p})$ and
\begin{equation*}
\|\chi_{(\nu/2,5\nu/4)}(y)w_{A,B}(y)^{-1}\Phi_\nu(y)\|_{p'}
\le C\nu^{-B+1-1/p}
\begin{cases}
\nu^{-1/2+2/(3p)}, & \text{ if } p<4/3,\\
(\log \nu)^{1/4}, & \text{ if } p=4/3,\\
1, & \text{ if } p>4/3.\\
\end{cases}
\end{equation*}
Now, considering separately the cases $p<4/3$, $p=4/3$, $p>4/3$ and using
appropriate assumptions on the relation between $b$ and $B$ proves the claim.

Consider the case~$i=3$. This time we bound $F_3$ by a constant times
\begin{equation*}
\|\chi_{(0,\nu/2)}(x)w_{a,b}(x)\Phi_\nu(x) \|_{p}\cdot
\|\chi_{(2\nu/3,\infty)}(y)y^{-(B+2)}\Phi_\nu(y)\|_{p'}\cdot
\|f(y)w_{A,B}(y)\|_p.
\end{equation*}
Similarly to (\ref{ec:estimacion0}) (the assumption
$a>-1/p-(\alpha+1/2)$, with $a\ge-(\alpha+1/2)$ when $p=\infty$, is used),
we obtain the estimate
\begin{equation*}
\|\chi_{(0,\nu/2)}(x)w_{a,b}(x)\Phi_\nu(x) \|_{p}
\le C 2^{-\nu}\Gamma(\nu+1)^{-1}(\nu/2)^{b+\nu+1/2}.
\end{equation*}
Moreover,
\begin{equation*}
\|\chi_{(2\nu/3,\infty)}(y)y^{-(B+2)}\Phi_\nu(y)\|_{p'} \le C
\nu^{-B-1-1/p}
\begin{cases}
\nu^{-1/2+2/(3p)}, & \text{ if } p<4/3,\\
(\log \nu)^{1/4}, & \text{ if } p=4/3,\\
1, & \text{ if } p>4/3.\\
\end{cases}
\end{equation*}
This is easily seen by splitting the integration onto $(2\nu/3,2\nu)$ and
$(2\nu,\infty)$, using the relevant assumption and then by inspection.
Combining these two bounds shows that the product of the first two
terms majorizing $F_3$ is a bounded function of $\nu\ge1$.

Finally, consider $i=4$. Then $F_4$ is bounded by a constant times
\begin{equation*}
\|\chi_{(\nu/2,5\nu/4)}(x)x^b\Phi_\nu(x)\|_{p}\cdot
\|\chi_{(3\nu/2,\infty)}(y)y^{-(B+2)}\|_{p'}\cdot
\|f(y)w_{A,B}(y)\|_p.
\end{equation*}
Here we use $\|\chi_{(3\nu/2,\infty)}(y)y^{-(B+2)}\|_{p'}=
O(\nu^{-B-1-1/p})$ (the assumption $B>-1-1/p$, $B\ge-2$ when $p=1$, is applied)
and (\ref{ec:normapesob}) with $b$ instead of $b-2$ (to be precise with
$(\nu/2,5\nu/4)$ in place of $(2\nu/3,\infty)$). Then,
considering separately the cases $p<4/3$, $p=4/3$, $p>4/3$ and using
appropriate assumptions on the relation between $b$ and $B$ shows that
the product of the two relevant terms is a bounded function of $\nu\ge1$.
\end{proof}

%%%%%%%%%%%%%%%%%%%%%%%%%%%%%%%%%%%%%%%%%%%%%%%%%%%%%%
\section{Proof of Theorem~\ref{Theorem:main}}
%%%%%%%%%%%%%%%%%%%%%%%%%%%%%%%%%%%%%%%%%%%%%%%%%%%%%%

We start with the proof of sufficiency in
Theorem~\ref{Theorem:main}. If $\alpha<1$ and $\nu\in[\alpha,1]$,
then (\ref{t4}) is a consequence of the estimates contained in
Propositions~\ref{Prop:P1}, \ref{Prop:P2} and~\ref{Prop:P4}. If
$\alpha\ge1$ and $\nu\in[\alpha,\infty)$, then we use the
estimates from Propositions \ref{Prop:P1}--\ref{Prop:P4} (note
that the regions that appear there cover the quarter plane
$(0,\infty)\times(0,\infty)$). Checking the required assumptions
shows that (\ref{t1}), (\ref{t2}) and (\ref{t3}) were used.

We now pass to the necessity. Checking of (\ref{t1}) relies on the
asymptotics of $K^1_\alpha(x,y)$ at $(0,0)$. Indeed,
to check the necessity of the first condition in (\ref{t1}) note that for a
sufficiently small $\varepsilon=\varepsilon(\alpha)$ we have $J_\alpha(s)\ge
Cs^\alpha$ for $0<s<\varepsilon$. Hence, for $x,y\in (0,\varepsilon)$,
$K^1_\alpha(x,y)\ge C(xy)^{\alpha+1/2}$. Take $f=\chi_{(\varepsilon/2,
\varepsilon)}$; then, for $x\in(0,\varepsilon)$,
$$
S_\alpha^1f(x)= \int_0^\infty K^1_\alpha(x,y) f(y)\,dy
\ge Cx^{\alpha+1/2} \int_{\varepsilon/2}^{\varepsilon}y^{\alpha+1/2}\,dy
\ge Cx^{\alpha+1/2}.
$$
Thus, (here and later on an appropriate interpretation is needed if
$p=\infty$), by (\ref{t4}),
\begin{align*}
\bigg(\int_0^\varepsilon x^{(\alpha+1/2+a)p}\,dx\bigg)^{1/p} &\le C
\bigg(\int_{0}^{\varepsilon}|S^1_\alpha f(x) x^a(1+x)^{b-a}|^p\,dx\bigg)^{1/p}\\
&\le C\bigg(\int_{\varepsilon/2}^{\varepsilon}| x^A(1+x)^{B-A}|^p\,dx\bigg)^{1/p}\\
&<\infty.
\end{align*}
In consequence, $a>-1/p-(\alpha+1/2)$ ($\ge$ if $p=\infty$) follows.

To verify the necessity of the second part of (\ref{t1}) define the function
$f$ by
$$
f(x) =
\begin{cases}
x^{-1-A}/({-\log x})^2, & \text{ if } p=1,\\
x^{-1/p-A}/({-\log x}), & \text{ if } 1<p<\infty,\\
x^{-A}, & \text{ if } p=\infty,
\end{cases}
$$
for $0<x<\varepsilon$ and $f(x)=0$ otherwise, and
consider in (\ref{t4}) the sequence of functions $f_n(x)=f(x)
\chi_{(\varepsilon/n,\varepsilon)}(x)$. Then,
$\sup_{n\ge1}\|f_n(x)x^A(1+x)^{B-A}\|_p<\infty$
and, for $0<x<\varepsilon$,
$$
S_\alpha^1f_n(x) = \int_0^\infty K^1_\alpha(x,y) f_n(y)\,dy
\ge Cx^{\alpha+1/2} \int_{\varepsilon/n}^{\varepsilon}y^{\alpha+1/2}f(y)\,dy.
$$
Thus, to keep $\|S_\alpha^1f_n(x)x^a(1+x)^{b-a}\|_p$ uniformly bounded
(with respect to $n$) requires $\alpha+1/2-A-1/p>-1$ ($\ge$ if $p=1$),
the second part of (\ref{t1}).

To check the necessity of the first condition in (\ref{t2}) choose a
sufficiently small $\varepsilon$ such that $J_\alpha(y)>0$ for $0<y<
\varepsilon$ and take $f=\chi_{(\varepsilon/2,\varepsilon)}$. Then, $S^1_\alpha
f(x)$ is the sum of
\begin{equation}
2\int_{\varepsilon/2}^\varepsilon \frac{2\alpha^2-(x^2+y^2)}{(x^2-y^2)^2}
\,y^{1/2}J_\alpha(y)\,dy \cdot x^{1/2}J_\alpha(x),
\label{A}
\end{equation}
\begin{equation}
-4x\int_{\varepsilon/2}^\varepsilon \frac{x^2+y^2}{(x^2-y^2)^3}
\,y^{1/2}J_\alpha(y)\,dy \cdot x^{1/2}J_\alpha'(x),
\label{B}
\end{equation}
\begin{equation}
4\int_{\varepsilon/2}^\varepsilon \frac{y(x^2+y^2)}{(x^2-y^2)^3}
\,y^{1/2}J_\alpha'(y)\,dy \cdot x^{1/2}J_\alpha(x),
\label{C}
\end{equation}
\begin{equation}
-4x\int_{\varepsilon/2}^\varepsilon \frac{y}{(x^2-y^2)^2}
\,y^{1/2}J_\alpha'(y)\,dy \cdot x^{1/2}J_\alpha'(x).
\label{D}
\end{equation}
Since
$$
x^{1/2}J_\alpha(x) = (2/\pi)^{1/2} \cos\Big(x-\frac{\alpha\pi}2-
\frac\pi4\Big) + O(x^{-1}), \qquad x\ge1,
$$
and
$$
J_\alpha'(x) = \frac{\alpha}xJ_\alpha(x)-J_{\alpha+1}(x),
$$
hence
$$
x^{1/2}J_\alpha'(x) = (2/\pi)^{1/2} \sin\Big(x-\frac{\alpha\pi}2-
\frac\pi4\Big) + O(x^{-1}), \qquad x\ge1.
$$
Let $g$ and $h$, $4<g<h<\infty$, be such that $x^{1/2}J_\alpha(x)\ge
1/\sqrt{2\pi}$ for $x\in X=\cup_{j=0}^\infty(g+2\pi j,h+2\pi j)$.
Then, for $x\in X$, the absolute values of (\ref{B}), (\ref{C}) and
(\ref{D}) are less than $Cx^{-3}$, while the (negative) value of (\ref{A})
is less than $-Cx^{-2}$. Thus, for large values of $x\in X$, the sum of absolute values
of (\ref{B}), (\ref{C}) and (\ref{D}) is less than the half of the absolute
value of~(\ref{A}). Therefore, it is possible to choose $g$ so large that the former
statement remains valid for all $x\in X$. Hence, with such a choice of $g$, for
$x\in X$, $|S^1_\alpha f(x)|\ge Cx^{-2}$, and using (\ref{t4}) produces
\begin{align*}
\bigg(\int_Xx^{(b-2)p}\,dx\bigg)^{1/p}&
\le C\bigg(\int_X|S^1_\alpha f(x)x^a(1+x)^{b-a}|^p\,dx\bigg)^{1/p}\\
&\le C\bigg(\int_{\varepsilon/2}^\varepsilon|f(x)x^A(1+x)^{B-A}|^p\,dx\bigg)^{1/p}\\
&<\infty.
\end{align*}
In consequence, $b<2-1/p$ ($\le$ if $p=\infty$) follows.

To check the necessity of the second condition in (\ref{t2}), define the
function $f$ by
$$
f(x) =
\begin{cases}
x^{-1-B}/(\log x)^{2}, & \text{ if } p=1,\\
x^{-1/p-B}/(\log x), & \text{ if } 1<p<\infty,\\
x^{-B}, & \text{ if } p=\infty,
\end{cases}
$$
for $x\ge2$ and $f(x)=0$ otherwise. Then choose an interval $(g,h)$,
$1<g<h<\infty$, such that $J_\alpha(x)<0$ and $J_\alpha'(x)>0$ for $x\in(g,h)$.
In addition, choose $G$ and $H$, $2\max\{g,\alpha\}<G<H<\infty$ in such a way
that $y^{1/2}J_\alpha(y)\ge 1/\sqrt{2\pi}$ and $J_\alpha'(y)\le 0$ for
$y\in Y=\cup_{j=0}^\infty(G+2\pi j,H+2\pi j)$.
Consider in (\ref{t4}) the sequence of functions $f_n(x)=f(x)\chi_{Y_n}(x)$,
$Y_n=Y\cap(2,n)$. Then, $\sup_{n\ge1}\|f_n(x)x^A(1+x)^{B-A}\|_p<\infty$
and $S^1_\alpha f_n(x)$ is the sum of
\begin{align*}
I_1(n,x) &= 2\int_{Y_n}f(y)\frac{2\alpha^2-(x^2+y^2)}{(x^2-y^2)^2}
\,y^{1/2}J_\alpha(y)\,dy\cdot x^{1/2}J_\alpha(x),\\
I_2(n,x) &= -4x\int_{Y_n}f(y)\frac{x^2+y^2}{(x^2-y^2)^3}
\,y^{1/2}J_\alpha(y)\,dy\cdot x^{1/2}J_\alpha'(x),\\
I_3(n,x) &= 4\int_{Y_n}f(y)\frac{y(x^2+y^2)}{(x^2-y^2)^3}
\,y^{1/2}J_\alpha'(y)\,dy\cdot x^{1/2}J_\alpha(x),\\
I_4(n,x) &= -4x\int_{Y_n}f(y)\frac{y}{(x^2-y^2)^2}
\,y^{1/2}J_\alpha'(y)\,dy\cdot x^{1/2}J_\alpha'(x).
\end{align*}
The choice of $f$, $(g,h)$ and $Y$ shows that for $x\in(g,h)$, $I_1(n,x)$,
$I_2(n,x)$ and $I_4(n,x)$ are positive while $I_3(n,x)$ in negative. Thus
$I_1(n,x)+I_3(n,x)\le S^1_\alpha f_n(x)$. Moreover,
$$
0 < \int_{Y_n}\frac{f(y)}{y^2}\,dy \cdot
\big({-x^{1/2}}J_\alpha(x)\big)\le C I_1(n,x)
$$
and
$$
|I_3(n,x)|\le C\int_{Y_n}\frac{f(y)}{y^3}\,dy \cdot
\big({-x^{1/2}}J_\alpha(x)\big).
$$
The last estimate shows that $G$ might be chosen so large that
$|I_3(n,x)|\le 10^{-2}I_1(n,x)$ for $x\in(g,h)$. With such a choice of $G$,
$$
\int_{Y_n}\frac{f(y)}{y^2}\,dy \cdot
\big({-x^{1/2}}J_\alpha(x)\big)\le CS^1_\alpha
f_n(x)
$$
and then multiplication of both sides by $x^b$ and $x$-integration
over $(g,h)$ produces
\begin{align*}
\int_{Y_n}\frac{f(y)}{y^2}\,dy
\cdot \bigg(\int_g^h\big|J_\alpha(x)x^{(b+1/2)} \big|^p\,dx\bigg)^{1/p}
&\le C\bigg(\int_g^h|S^1_\alpha f_n(y)y^b|^p\,dy\bigg)^{1/p}\\
&\le C\|f_n(x)x^A(1+x)^{B-A}\|_p.
\end{align*}
Thus, uniform boundedness (with respect to $n$) of the integrals
$\int_{Y_n}f(y)y^{-2}\,dy$ now requires $B>-1-1/p$ ($\ge$ if $p=1$),
the second part of (\ref{t2}).

Proving the necessity of (\ref{t3}) note that for $R>0$,
$K_{\alpha,R}^1(x,y)=R K_{\alpha,1}^1(xR,yR)$, hence
$S_{\alpha,R}^1f(x)= S_{\alpha,1}^1f_R(xR)$,
where $f_R(y)=f(y/R)$. Then, applying (\ref{t4}) (where $R=1$ is assumed)
to a nontrivial function $f\in C^\infty_c(0,\infty)$ produces
\begin{align}
\|S_{\alpha,R}^1 f(x)x^a(R^{-1}+x)^{b-a}\|_p
&\le C R^{B-b}\|f(x)x^A(R^{-1}+x)^{B-A})\|_p \nonumber\\
&\le C R^{B-b}\|f(x)x^B\|_p
\label{n1}
\end{align}
if $B-A\le0$ or with the last term replaced by $\|f(x)x^A(1+x)^{B-A}\|_p$ if
$B-A>0$ and $R>1$. But
\begin{equation}
\lim_{R\to\infty}S_{\alpha,R}^1 f(x)=f(x), \qquad x>0,
\label{n2}
\end{equation}
hence by (\ref{n1}), (\ref{n2}) and Fatou's lemma for $p<\infty$ or trivially for
$p=\infty$ we have
\begin{equation}
\|f(x)x^b\|_p\le C \|f(x)x^B\|_p\liminf_{R\to\infty}R^{B-b}
\end{equation}
if $B-A\le0$ or analogous inequality with $x^B$ replaced by $x^A(1+x)^{B-A}$
when $B-A>0$. In any case, this implies $b\le B$. To verify (\ref{n2}),
fix $x>0$ and $f\in C^\infty_c(0,\infty)$ and use $S_{\nu,R}^1f(x)
= {\mathcal H}_\nu(m_R^1\cdot{\mathcal H}_\nu f)(x)$ to write
\begin{equation*}
S_{\nu,R}^1 f(x) = \int_0^\infty (xy)^{1/2}J_\alpha(xy)\Big(1-\Big(
\frac yR\Big)^2\Big)_+{\mathcal H}_\alpha f(y)\,dy.
\end{equation*}
Then, (\ref{n2}) follows by an application of the dominated
convergence theorem and the inversion formula for the Hankel
transform ${\mathcal H}_\alpha$ once we note that
$$
\int_0^\infty (xy)^{1/2}|J_\alpha(xy)||{\mathcal H}_\alpha f(y)|\,dy
< \infty.
$$
The above is a consequence of ${\mathcal H}_\alpha
f(y)=O(y^{\alpha+1/2})$, $y\to0^+$, which in turn follows from
(\ref{asy}) and ${\mathcal H}_\alpha f(y) = O(y^{-2})$, the latter
easily implied by (\ref{asy2}). This finishes the proof of the
necessity of (\ref{t3}), hence the theorem.

%%%%%%%%%%%%%%%%%%%%%%%%%%%%%%%%%%%%%%%%%%%%%%%%%%%%%%
\begin{thebibliography}{0}
%%%%%%%%%%%%%%%%%%%%%%%%%%%%%%%%%%%%%%%%%%%%%%%%%%%%%%

\bibitem{BC} J. A. Barcel\'o and A. C\'ordoba,
Band-limited functions: $L^p$-convergence,
\textit{Trans. Amer. Math. Soc.}
\textbf{313} (1989), 655--669.

\bibitem{C} A. C\'ordoba, The disc multiplier,
\textit{Duke Math. J.} \textbf{58} (1989), 21--27.

\bibitem{Le} N. N. Lebedev,
Special functions and its applications,
Dover, New York, 1972.

\bibitem{Mu} B. Muckenhoupt,
Mean convergence of Hermite and Laguerre series. I,
\textit{Trans. Amer. Math. Soc.}
\textbf{147} (1970), 419--431;
{II}, \textit{ibid.} \textbf{147} (1970), 433--460.

\bibitem{Mu-We} B. Muckenhoupt and D. W. Webb,
Two-weight norm inequalities for Ces\`aro means of Laguerre expansions,
\textit{Trans. Amer. Math. Soc.}
\textbf{353} (2000), 1119--1149.

\bibitem{Po} E. L. Poiani,
Mean Ces\`aro summability of Laguerre and Hermite series,
\textit{Trans. Amer. Math. Soc.}
\textbf{173} (1972), 1--31.

\bibitem{St} K. Stempak,
Uniform two-weight norm inequalities for Hankel transform partial
sum operators,
\textit{Houston J. Math.}
\textbf{29} (2003), 1045--1063.

\bibitem{We} D. W. Webb,
Pointwise estimates of the moduli of Ces\`aro-Laguerre kernels,
{preprint}.

\end{thebibliography}

\end{document}

%%%%%%%%%%%%%%%%%%%%%%%%%%%%%%%%


